% GENERATED FILE. Edit canonical lessons, then run npm run generate:books.
% canonical-source-hash: fnv1a64:ee803aba9a9fb3ac
% canonical-lessons: MW-C272-phagat, MW-C272-bi, MW-C272-jarur, MW-C272-besak, MW-C272-manjur, MW-R272-review-266-272

\chapter{Only This, and Me Too}
\label{ch:mw-only-this-and-me-too}
\hlchaptermodality{272}

\begin{chapteropening}
\textbf{By the end of this chapter:} \emph{I can say only, also, of course and certainly.}

Complete the last lesson of chapter 272: I can say only, also, of course and certainly.
\end{chapteropening}

\section[\texorpdfstring{\mw{फगत} (phagat)}{फगत (phagat)}]{\mw{फगत} (phagat) — only}

\label{lesson:MW-C272-phagat}



\begin{quote}
Before the new one: say the Marwadi for a witness, then the Marwadi for an oath.
\end{quote}

\subsection*{You'll want to know: \mw{फगत}}
\textbf{\mw{फगत}} — \emph{phagat} — \textquotedblleft{}only\textquotedblright{}.

\begin{quote}
\mw{हूं} \mw{फगत} \mw{चाय} \mw{पीऊं} \mw{हूं।}
\end{quote}

\begin{quote}
\emph{hū̃ phagat chāy pīū̃ hū̃.}
\end{quote}

\begin{quote}
I only drink tea.
\end{quote}

\begin{cousinweb}
\textbf{\mw{फगत}} (\emph{phagat}) is usually traced, through Persian \emph{faqat}, to Arabic \emph{faqaṭ}, \textquotedblleft{}only\textquotedblright{}. It goes right before the thing it limits.
\end{cousinweb}

\begin{grammarlens}[title={What you've built}]
One more word for asking and answering.
\end{grammarlens}

\subsection*{Your turn}
\begin{itemize}
\raggedright
  \item \emph{Say it:} \emph{phagat}
  \item \emph{Say it:} \emph{phagat}, once more
  \item \emph{Say it:} \emph{hū̃ phagat chāy pīū̃ hū̃.} — I only drink tea.
  \item \emph{Recall:} say the Marwadi for a witness, then the Marwadi for an oath, then say \emph{phagat} again
\end{itemize}

\begin{tcolorbox}[breakable,colback=teal!4,colframe=teal!35!black,title={Before you move on}]
What does \mw{फगत} mean? (\textquotedblleft{}Only\textquotedblright{}.) Say it once more.
\end{tcolorbox}
\section[\texorpdfstring{\mw{बी} (bī)}{बी (bī)}]{\mw{बी} (bī) — also}

\label{lesson:MW-C272-bi}



\begin{quote}
Before the new one: say the Marwadi for an oath, then the Marwadi for only.
\end{quote}

\subsection*{You'll want to know: \mw{बी}}
\textbf{\mw{बी}} — \emph{bī} — \textquotedblleft{}also\textquotedblright{}.

\begin{quote}
\mw{हूं} \mw{बी} \mw{आसूं।}
\end{quote}

\begin{quote}
\emph{hū̃ bī āsū̃.}
\end{quote}

\begin{quote}
I'll come too.
\end{quote}

\begin{cousinweb}
\textbf{\mw{बी}} (\emph{bī}) is inherited from Sanskrit \emph{api ca}, \textquotedblleft{}and also\textquotedblright{}; CDIAL lists Old Marwari \emph{bī}, \textquotedblleft{}also\textquotedblright{} (CDIAL 486). It follows the word it adds.
\end{cousinweb}

\begin{grammarlens}[title={What you've built}]
One more word for asking and answering.
\end{grammarlens}

\subsection*{Your turn}
\begin{itemize}
\raggedright
  \item \emph{Say it:} \emph{bī}
  \item \emph{Say it:} \emph{bī}, once more
  \item \emph{Say it:} \emph{hū̃ bī āsū̃.} — I'll come too.
  \item \emph{Recall:} say the Marwadi for an oath, then the Marwadi for only, then say \emph{bī} again
\end{itemize}

\begin{tcolorbox}[breakable,colback=teal!4,colframe=teal!35!black,title={Before you move on}]
What does \mw{बी} mean? (\textquotedblleft{}Also\textquotedblright{}.) Say it once more.
\end{tcolorbox}
\section[\texorpdfstring{\mw{जरूर} (jarūr)}{जरूर (jarūr)}]{\mw{जरूर} (jarūr) — certainly}

\label{lesson:MW-C272-jarur}



\begin{quote}
Before the new one: say the Marwadi for only, then the Marwadi for also.
\end{quote}

\subsection*{You'll want to know: \mw{जरूर}}
\textbf{\mw{जरूर}} — \emph{jarūr} — \textquotedblleft{}certainly\textquotedblright{}.

\begin{quote}
\mw{हूं} \mw{जरूर} \mw{आसूं।}
\end{quote}

\begin{quote}
\emph{hū̃ jarūr āsū̃.}
\end{quote}

\begin{quote}
I'll certainly come.
\end{quote}

\begin{cousinweb}
\textbf{\mw{जरूर}} (\emph{jarūr}) is usually traced, through Persian \emph{zarūr}, \textquotedblleft{}necessary; certainly\textquotedblright{}, to Arabic \emph{ḍarūra}, \textquotedblleft{}necessity\textquotedblright{}. \textbf{\mw{जरूरी}}, necessary, and \textbf{\mw{जरूरत}}, a need, come from the same word.
\end{cousinweb}

\begin{grammarlens}[title={What you've built}]
One more word for asking and answering.
\end{grammarlens}

\subsection*{Your turn}
\begin{itemize}
\raggedright
  \item \emph{Say it:} \emph{jarūr}
  \item \emph{Say it:} \emph{jarūr}, once more
  \item \emph{Say it:} \emph{hū̃ jarūr āsū̃.} — I'll certainly come.
  \item \emph{Recall:} say the Marwadi for only, then the Marwadi for also, then say \emph{jarūr} again
\end{itemize}

\begin{tcolorbox}[breakable,colback=teal!4,colframe=teal!35!black,title={Before you move on}]
What does \mw{जरूर} mean? (\textquotedblleft{}Certainly\textquotedblright{}.) Say it once more.
\end{tcolorbox}
\section[\texorpdfstring{\mw{बेसक} (besak)}{बेसक (besak)}]{\mw{बेसक} (besak) — of course}

\label{lesson:MW-C272-besak}



\begin{quote}
Before the new one: say the Marwadi for also, then the Marwadi for certainly.
\end{quote}

\subsection*{You'll want to know: \mw{बेसक}}
\textbf{\mw{बेसक}} — \emph{besak} — \textquotedblleft{}of course\textquotedblright{}.

\begin{quote}
\mw{बेसक}, \mw{हूं} \mw{आसूं।}
\end{quote}

\begin{quote}
\emph{besak, hū̃ āsū̃.}
\end{quote}

\begin{quote}
Of course I'll come.
\end{quote}

\begin{cousinweb}
\textbf{\mw{बेसक}} (\emph{besak}) is usually traced to Persian \emph{bē-shakk}, \textquotedblleft{}without doubt\textquotedblright{}, built from Persian \emph{bē}, \textquotedblleft{}without\textquotedblright{}, and \emph{shakk}, \textquotedblleft{}doubt\textquotedblright{}, which Persian took from Arabic.
\end{cousinweb}

\begin{grammarlens}[title={What you've built}]
One more word for asking and answering.
\end{grammarlens}

\subsection*{Your turn}
\begin{itemize}
\raggedright
  \item \emph{Say it:} \emph{besak}
  \item \emph{Say it:} \emph{besak}, once more
  \item \emph{Say it:} \emph{besak, hū̃ āsū̃.} — Of course I'll come.
  \item \emph{Recall:} say the Marwadi for also, then the Marwadi for certainly, then say \emph{besak} again
\end{itemize}

\begin{tcolorbox}[breakable,colback=teal!4,colframe=teal!35!black,title={Before you move on}]
What does \mw{बेसक} mean? (\textquotedblleft{}Of course\textquotedblright{}.) Say it once more.
\end{tcolorbox}
\section[\texorpdfstring{\mw{मंजूर} (mañjūr)}{मंजूर (mañjūr)}]{\mw{मंजूर} (mañjūr) — agreed}

\label{lesson:MW-C272-manjur}



\begin{quote}
Before the new one: say the Marwadi for certainly, then the Marwadi for of course.
\end{quote}

\subsection*{You'll want to know: \mw{मंजूर}}
\textbf{\mw{मंजूर}} — \emph{mañjūr} — \textquotedblleft{}agreed\textquotedblright{}.

\begin{quote}
\mw{यो} \mw{भाव} \mw{मंजूर} \mw{है।}
\end{quote}

\begin{quote}
\emph{yo bhāv mañjūr hai.}
\end{quote}

\begin{quote}
This price is agreed.
\end{quote}

\begin{cousinweb}
\textbf{\mw{मंजूर}} (\emph{mañjūr}) is usually traced, through Persian \emph{manzūr}, to Arabic \emph{manẓūr}, \textquotedblleft{}looked at\textquotedblright{}, hence \textquotedblleft{}approved\textquotedblright{}.
\end{cousinweb}

\begin{grammarlens}[title={What you've built}]
One more word for asking and answering.
\end{grammarlens}

\subsection*{Your turn}
\begin{itemize}
\raggedright
  \item \emph{Say it:} \emph{mañjūr}
  \item \emph{Say it:} \emph{mañjūr}, once more
  \item \emph{Say it:} \emph{yo bhāv mañjūr hai.} — This price is agreed.
  \item \emph{Recall:} say the Marwadi for certainly, then the Marwadi for of course, then say \emph{mañjūr} again
\end{itemize}

\begin{tcolorbox}[breakable,colback=teal!4,colframe=teal!35!black,title={Before you move on}]
What does \mw{मंजूर} mean? (\textquotedblleft{}Agreed\textquotedblright{}.) Say it once more.
\end{tcolorbox}
\section[(dialogue)]{Looking back: from Please Come In to Only This, and Me Too}

\label{lesson:MW-R272-review-266-272}



\begin{quote}
Say the Marwadi for at your service and help.
\end{quote}

\subsection*{Your turn}
Say these aloud:

\begin{itemize}
\raggedright
  \item at your service, help, a kindness, a petition, a blessing
  \item permission, leave to go, an order, waiting, to hold on
  \item to say, to recite, counsel, a meaning, a good word for someone
  \item an excuse, a refusal, to dig in one's heels, stubbornness, to interrupt
  \item pardon, a wrongdoing, to regret, to smash, to get cross
  \item correct, wrong, proof, a witness, an oath
  \item only, also, certainly, of course, agreed
\end{itemize}

\begin{tcolorbox}[breakable,colback=teal!4,colframe=teal!35!black,title={Before you move on}]
At your service? (\textbf{\mw{हाजर}}, \emph{hājar}.) Help? (\textbf{\mw{मदद}}, \emph{madad}.) A kindness? (\textbf{\mw{मेहरबानी}}, \emph{meharbānī}.) A petition? (\textbf{\mw{अरज}}, \emph{araj}.) A blessing? (\textbf{\mw{आसीस}}, \emph{āsīs}.) Permission? (\textbf{\mw{इजाजत}}, \emph{ijājat}.) Leave to go? (\textbf{\mw{रजा}}, \emph{rajā}.) An order? (\textbf{\mw{हुकम}}, \emph{hukam}.) Waiting? (\textbf{\mw{इंतजार}}, \emph{intjār}.) To hold on? (\textbf{\mw{थमणो}}, \emph{thamṇo}.) To say? (\textbf{\mw{कैवणो}}, \emph{kaivṇo}.) To recite? (\textbf{\mw{सुणावणो}}, \emph{suṇāvṇo}.) Counsel? (\textbf{\mw{सलाह}}, \emph{salāh}.) A meaning? (\textbf{\mw{मतलब}}, \emph{matlab}.) A good word for someone? (\textbf{\mw{सिफारस}}, \emph{siphāras}.) An excuse? (\textbf{\mw{बहानो}}, \emph{bahāno}.) A refusal? (\textbf{\mw{इनकार}}, \emph{inkār}.) To dig in one's heels? (\textbf{\mw{अड़णो}}, \emph{aṛṇo}.) Stubbornness? (\textbf{\mw{जिद}}, \emph{jid}.) To interrupt? (\textbf{\mw{टोकणो}}, \emph{ṭokṇo}.) Pardon? (\textbf{\mw{माफी}}, \emph{māphī}.) A wrongdoing? (\textbf{\mw{कसूर}}, \emph{kasūr}.) To regret? (\textbf{\mw{पछताणो}}, \emph{pachhtāṇo}.) To smash? (\textbf{\mw{फोड़णो}}, \emph{phoṛṇo}.) To get cross? (\textbf{\mw{खीजणो}}, \emph{khījṇo}.) Correct? (\textbf{\mw{खरो}}, \emph{kharo}.) Wrong? (\textbf{\mw{गलत}}, \emph{galat}.) Proof? (\textbf{\mw{सबूत}}, \emph{sabūt}.) A witness? (\textbf{\mw{साखी}}, \emph{sākhī}.) An oath? (\textbf{\mw{कसम}}, \emph{kasam}.) Only? (\textbf{\mw{फगत}}, \emph{phagat}.) Also? (\textbf{\mw{बी}}, \emph{bī}.) Certainly? (\textbf{\mw{जरूर}}, \emph{jarūr}.) Of course? (\textbf{\mw{बेसक}}, \emph{besak}.) Agreed? (\textbf{\mw{मंजूर}}, \emph{mañjūr}.)
\end{tcolorbox}
